\documentclass[10pt,twocolumn]{article}

\usepackage[letterpaper,margin=0.65in,columnsep=0.25in]{geometry}
\usepackage[T1]{fontenc}
\usepackage[utf8]{inputenc}
\usepackage{amsmath}
\usepackage{amssymb}
\usepackage{mathtools}
\usepackage{enumitem}
\usepackage{microtype}
\usepackage{xcolor}
\usepackage{booktabs}
\usepackage{array}
\usepackage{graphicx}
\usepackage[hidelinks]{hyperref}

\setlength{\parindent}{0pt}
\setlength{\parskip}{2pt}
\setlist{nosep,leftmargin=*}
\renewcommand{\arraystretch}{1.08}
\microtypesetup{expansion=false}

\newcommand{\History}{\mathcal{H}}
\newcommand{\Frontiers}{\mathcal{F}}
\newcommand{\Memory}{\mathcal{M}}
\newcommand{\Global}{\mathcal{G}}
\newcommand{\Reveal}{\mathcal{Y}}
\newcommand{\Option}{\boldsymbol{\omega}}
\newcommand{\Loss}{\mathcal{L}}
\newcommand{\Expect}{\mathbb{E}}
\newcommand{\KL}{D_{\mathrm{KL}}}

\title{\vspace{-1.5cm}
\textbf{FrontierWorld: Counterfactual Revelation Models and\\
Lineage-Aware Predictive Memory for Object Navigation}
\vspace{-0.35cm}
}

\author{
Ashrith Edukulla \qquad Behrad Rabiei \qquad Yulun Tian\\
University of Michigan, Ann Arbor
}

\date{}

\begin{document}

\maketitle
\vspace{-0.75cm}

\begin{abstract}
Object navigation requires an agent to choose among multiple boundaries of its
known world while reasoning about what each choice may reveal. Existing
frontier methods largely treat a frontier as a temporary waypoint to score,
while predictive navigation methods complete a global map or generate a single
future without preserving a decision-aligned memory for each evolving
frontier. We propose \textit{FrontierWorld}, a framework that represents every
frontier as a persistent probabilistic belief over the information that would
be revealed by crossing it. The framework has three central components.
First, a frontier lineage graph preserves predictive state through frontier
birth, update, split, merge, and retirement. Second, a fixed-budget learned
memory retains observations according to their value for predicting future
frontier revelations, rather than according to recency alone. Third, an
action-conditioned stochastic world model predicts the geometry, semantics,
target evidence, traversability, and new exploration opportunities that would
be observed under each candidate frontier-crossing option. To train and
evaluate this model without policy-selection bias, we introduce a
branch-complete data-generation protocol that forks an HM3D simulator state
and executes every available frontier option from the same navigation state.
This permits direct evaluation of prediction calibration, frontier-ranking
regret, memory efficiency, and closed-loop ObjectNav performance.
\end{abstract}

\vspace{-0.15cm}

\section{Motivation}

A frontier is the boundary between explored free space and an unresolved
region. Classical exploration methods treat it as a geometric waypoint
\cite{yamauchi1997}; recent systems enrich it with learned information gain
\cite{frontiernet}, visual-language relevance \cite{vlfm,openfrontier}, or
object and trajectory features \cite{fomnav}. Other methods infer the geometry
and semantics of unseen space through map completion \cite{foresightnav} or
predict future egocentric observations along actions \cite{nwm}.

These advances leave a different question unresolved:

\begin{quote}
\textit{What is the smallest frontier-specific state that an agent should
remember in order to predict the information that crossing that frontier would
reveal?}
\end{quote}

This question is important because a frontier is not a stable object. As new
observations arrive, a boundary can move, split, merge, or disappear. Moreover,
the relevant future is not simply a plausible image behind the boundary. The
agent needs a calibrated distribution over future observations and map
updates under a concrete action, together with the expected utility and risk
of acquiring them.

We therefore formulate frontier selection as counterfactual prediction under
partial observability. The key distinction is that FrontierWorld models
\emph{future information acquisition}: how the agent's belief will change
when it executes a frontier-crossing option.

\section{Problem Formulation}

At time \(t\), the agent has observation-action history

\begin{equation}
\History_t =
\left(o_1,a_1,\ldots,o_t\right),
\qquad
o_t=(I_t,D_t,P_t),
\end{equation}

where \(I_t\), \(D_t\), and \(P_t\) denote RGB, depth, and pose. A geometric or
visual frontier detector produces

\begin{equation}
\Frontiers_t=\{f_1^t,\ldots,f_{N_t}^t\}.
\end{equation}

For each candidate \(f_i^t\), a planner constructs a macro-action

\begin{equation}
\Option_i^t =
\left(\tau_i^{\mathrm{approach}},
\tau_i^{\mathrm{cross}},
H\right),
\end{equation}

consisting of a path to the frontier, a short probing trajectory into unknown
space, and a prediction horizon \(H\).

We define the \emph{frontier revelation} as

\begin{equation}
\Reveal_{t,i}^{H}
=
\operatorname{Reveal}
\left(
s,o_{t+1:t+H}
\mid
\operatorname{do}(\Option_i^t)
\right),
\label{eq:revelation}
\end{equation}

where \(s\) is the underlying scene. The intervention notation emphasizes that
all candidate options begin from the same state. The revelation contains

\begin{equation}
\Reveal_{t,i}^{H}
=
\left(
\Delta M_i^{\mathrm{occ}},
\Delta M_i^{\mathrm{sem}},
r_i,
y_i^g,
\chi_i,
\Frontiers_i^{\mathrm{new}}
\right),
\end{equation}

where \(\Delta M_i^{\mathrm{occ}}\) and
\(\Delta M_i^{\mathrm{sem}}\) are newly revealed occupancy and semantic
regions, \(r_i\) is room type, \(y_i^g\) indicates target evidence,
\(\chi_i\) is traversability or collision outcome, and
\(\Frontiers_i^{\mathrm{new}}\) is the set of newly exposed frontiers.

\section{Frontier Lineage Graph}

Nearest-neighbor matching assumes one frontier at time \(t\) corresponds to
one frontier at time \(t+1\). This fails when map revelation changes the
boundary topology. We instead maintain a lineage graph

\begin{equation}
\mathcal{L}_t=(V_t,E_t),
\end{equation}

whose nodes are time-indexed frontier observations and whose edges represent

\begin{equation}
e\in
\{
\mathrm{birth},
\mathrm{update},
\mathrm{split},
\mathrm{merge},
\mathrm{retire}
\}.
\end{equation}

Association uses boundary overlap in a common map frame, orientation,
visibility-ray agreement, appearance, and predicted connection to the same
unknown region. Let \(A_{ij}\in[0,1]\) be the soft lineage weight from
\(f_i^t\) to \(f_j^{t+1}\). A child frontier inherits aligned parent memory:

\begin{equation}
\widetilde m_j^{t+1}
=
\sum_{i}
A_{ij}
T_{i\rightarrow j}m_i^t,
\label{eq:inherit}
\end{equation}

where \(T_{i\rightarrow j}\) transforms spatial tokens into the child's
frontier-centered coordinate frame. Equation~\eqref{eq:inherit} naturally
supports one-to-many splits and many-to-one merges.

The graph distinguishes a physical unresolved region from the boundary
segments through which it is currently visible. Consequently, multiple
frontiers believed to enter the same latent room can share evidence without
being collapsed into one waypoint.

\section{Budgeted Predictive Frontier Memory}

Rather than manually storing every image or reusing architecture-dependent raw
KV entries, each frontier maintains \(B\) learned memory slots:

\begin{equation}
m_i^t =
W_{\psi}
\left(
\widetilde m_i^t,
\left\{
\Pi_i E(o_\tau)
\right\}_{\tau\leq t}
\right),
\qquad
\lvert m_i^t\rvert\leq B.
\end{equation}

Here \(E\) is a frozen or lightly adapted visual encoder and \(\Pi_i\) assigns
pose-grounded observation tokens to frontier \(i\). The writer
\(W_\psi\) performs fusion, compression, and eviction.

We train the memory as a predictive information bottleneck:

\begin{equation}
\max_{\psi}
\;
I
\left(
m_i^t;\Reveal_{t,i}^{H}
\mid
\Global_t,\Option_i^t
\right)
-
\beta
I
\left(
m_i^t;\History_t
\right),
\label{eq:ib}
\end{equation}

subject to the slot budget. The first term preserves history that predicts
future revelation; the second discourages indiscriminate memorization.
Operationally, Eq.~\eqref{eq:ib} is approximated by multi-task revelation
losses and a capacity constraint.

The value of a stored token \(x\) can be estimated by its effect on predictive
loss:

\begin{equation}
v(x)
\approx
\Loss_{\mathrm{pred}}
\left(m_i^t\setminus x\right)
-
\Loss_{\mathrm{pred}}
\left(m_i^t\right).
\end{equation}

This provides a learned eviction criterion and permits direct comparison with
FIFO, most-recent, visual-diversity, and uncertainty-based cache policies.

\section{Counterfactual Revelation World Model}

FrontierWorld predicts a distribution rather than one deterministic future:

\begin{equation}
p_\theta
\left(
\Reveal_{t,i}^{H}
\mid
m_i^t,\Global_t,\Option_i^t,g
\right)
=
\sum_{k=1}^{K}
\pi_{ik}
p_{\theta,k}
\left(
\Reveal_{t,i}^{H}
\mid
m_i^t,\Global_t,\Option_i^t,g
\right).
\label{eq:mixture}
\end{equation}

The \(K\) components represent distinct spatial-semantic hypotheses, such as a
doorway leading to a bedroom, hallway, bathroom, or closet. Predicting multiple
hypotheses avoids penalizing a calibrated model for representing genuine
ambiguity in unobserved space.

A practical first implementation uses frontier-centered RGB-D tokens, global
map tokens, relative-pose tokens describing \(\Option_i^t\), and a compact
Transformer or Perceiver decoder. The output heads predict:

\begin{itemize}
    \item local occupancy logits in a frontier-centered bird's-eye view;
    \item open-vocabulary semantic features or category logits;
    \item room and goal-presence distributions;
    \item newly visible area and newly created frontiers;
    \item traversability, collision risk, and uncertainty.
\end{itemize}

Future RGB or video is optional. The core model operates in decision-relevant
geometry and semantic feature space, avoiding the cost and evaluation
ambiguity of photorealistic generation.

The training objective is

\begin{align}
\Loss_{\mathrm{pred}}
={}&
\lambda_{\mathrm{occ}}\Loss_{\mathrm{occ}}
+
\lambda_{\mathrm{sem}}\Loss_{\mathrm{sem}}
+
\lambda_{\mathrm{goal}}\Loss_{\mathrm{goal}}
\nonumber\\
&+
\lambda_{\mathrm{trav}}\Loss_{\mathrm{trav}}
+
\lambda_{\mathrm{area}}\Loss_{\mathrm{area}}
+
\lambda_{\mathrm{cal}}\Loss_{\mathrm{cal}},
\end{align}

where \(\Loss_{\mathrm{cal}}\) is a proper probabilistic loss such as negative
log-likelihood or Brier loss.

\section{Belief-Consistent Frontier Planning}

For every option, the model estimates goal utility, expected information,
travel cost, and risk. The next frontier is

\begin{equation}
\begin{aligned}
f_t^\star
=
\arg\max_{f_i\in\Frontiers_t}
\big[
&\Expect_{\,\Reveal_i}U(\Reveal_i,g)
+\eta\,\mathcal{I}_i\\
&-\lambda C(\Option_i)
-\rho\,\operatorname{CVaR}_{\alpha}(R_i)
\big],
\end{aligned}
\label{eq:planner}
\end{equation}

where \(\mathcal{I}_i\) is expected information gain and
\(\operatorname{CVaR}_{\alpha}\) penalizes high-risk branches. A VLM may provide
semantic priors, but the final score in Eq.~\eqref{eq:planner} is grounded in
predicted outcomes rather than an opaque prompt response.

After executing \(f_t^\star\), the realized revelation
\(\Reveal_{t,\star}^{\mathrm{obs}}\) corrects its belief:

\begin{equation}
q_{t+1}(z_\star)
\propto
p
\left(
\Reveal_{t,\star}^{\mathrm{obs}}
\mid z_\star
\right)
q_t(z_\star).
\end{equation}

Prediction residuals are written to the frontier lineage and propagated to
frontiers believed to border the same unknown region. Thus, discovering that
one entrance leads to a target-absent room can reduce the target probability
of a second entrance into that room.

\section{Branch-Complete FrontierReveal Benchmark}

Standard navigation trajectories observe only the outcome of the selected
frontier. This prevents supervised comparison of all candidate choices and
introduces policy-selection bias. We propose a branch-complete simulator
protocol:

\begin{enumerate}
    \item run an exploration policy to a frontier decision state;
    \item save the simulator and agent memory state;
    \item enumerate all valid frontiers \(f_i\in\Frontiers_t\);
    \item fork the identical state for each \(f_i\);
    \item execute the canonical option \(\Option_i^t\);
    \item record the resulting revelation \(\Reveal_{t,i}^{H}\);
    \item restore the original state before evaluating the next branch.
\end{enumerate}

Each training example is

\begin{equation}
\mathcal{D}_t
=
\left\{
\left(
\History_t,\Global_t,f_i,\Option_i^t,
\Reveal_{t,i}^{H}
\right)
\right\}_{i=1}^{N_t}.
\end{equation}

This dataset permits counterfactual evaluation from the same partial
observation state. It also supplies ground-truth frontier genealogies by
tracking boundary overlap and revealed free-space connectivity in the
simulator.

\section{Experimental Plan}

\textbf{Environment.}
We use HM3D training scenes for data generation and held-out HM3D-v2 scenes
for evaluation. Candidate histories are collected with geometric,
information-gain, and semantically guided exploration policies to avoid
overfitting to one behavior distribution.

\textbf{Prediction metrics.}
We report occupancy IoU, semantic mIoU or feature similarity, room
classification, goal-presence AUROC and Brier score, traversability accuracy,
newly revealed area error, expected calibration error, and best-of-\(K\)
hypothesis coverage.

\textbf{Decision metrics.}
For predicted utility \(\widehat U\), frontier-selection regret is

\begin{equation}
\operatorname{Regret}_t
=
U
\left(
f_t^{\mathrm{oracle}}
\right)
-
U
\left(
f_t^{\mathrm{pred}}
\right).
\end{equation}

We additionally report Spearman correlation between predicted and oracle
frontier rankings and top-\(k\) oracle recall.

\textbf{Memory metrics.}
At budgets \(B\in\{0,1,4,8,16,32\}\), we measure prediction quality,
selection regret, memory usage, visual-token count, inference latency, and
closed-loop navigation.

\textbf{Lineage metrics.}
We report association precision and recall, IDF1, ID switches, cache
contamination, and prediction quality separately for stable, split, merge,
and long-absence frontier events.

\textbf{Navigation metrics.}
The final system is evaluated using Success Rate, SPL, SoftSPL, distance to
success, collisions, frontier revisits, oscillation rate, model calls, and
wall-clock inference cost.

\begin{table}[t]
\centering
\footnotesize
\setlength{\tabcolsep}{3.5pt}
\caption{Core ablations. Entries are reported at a common decision state and
memory budget; no-cache and shuffled-cache conditions control for scene priors
and additional token count.}
\label{tab:ablations}
\begin{tabular}{lcccc}
\toprule
Method & Pred. & Cal. & Regret & SPL\\
& \(\uparrow\) & \(\downarrow\) & \(\downarrow\) & \(\uparrow\)\\
\midrule
Current observation only & -- & -- & -- & --\\
Single frontier keyframe & -- & -- & -- & --\\
Global temporal memory & -- & -- & -- & --\\
Frontier FIFO memory & -- & -- & -- & --\\
Shuffled frontier memory & -- & -- & -- & --\\
Ours: learned frontier memory & -- & -- & -- & --\\
\bottomrule
\end{tabular}
\end{table}

\begin{table}[t]
\centering
\footnotesize
\setlength{\tabcolsep}{4pt}
\caption{Frontier persistence ablations. Oracle lineage provides an upper
bound; the shuffled condition measures sensitivity to memory contamination.}
\label{tab:lineage}
\begin{tabular}{lccc}
\toprule
Association & IDF1 \(\uparrow\) & Pred. \(\uparrow\) & Regret \(\downarrow\)\\
\midrule
No persistence & -- & -- & --\\
Nearest position & -- & -- & --\\
Position + appearance & -- & -- & --\\
Ours: lineage graph & -- & -- & --\\
Oracle lineage & -- & -- & --\\
\bottomrule
\end{tabular}
\end{table}

\section{Relation to Prior Work and Claimed Novelty}

Frontier-enhanced Topological Memory predicts features at unvisited but
already visible ghost-node viewpoints \cite{ftm}; FrontierWorld instead
predicts the distribution of \emph{new information} revealed after executing a
frontier-crossing option. ForesightNav completes a global geometric-semantic
map \cite{foresightnav}; FrontierWorld predicts stochastic, decision-indexed
branches and evaluates them against branch-complete counterfactual outcomes.
FOM-Nav encodes hand-designed frontier, object, and path features
\cite{fomnav}; FrontierWorld learns a bounded frontier memory specifically for
future-revelation prediction. Navigation World Models generate
action-conditioned future views \cite{nwm}; FrontierWorld contributes an
evolving frontier address space, lineage-preserving memory, and explicit
uncertainty over information acquisition.

The intended contributions are:

\begin{enumerate}
    \item a formulation of frontier selection as action-conditioned prediction
    of future information acquisition;
    \item a lineage graph that preserves predictive belief through frontier
    birth, split, merge, and retirement;
    \item a fixed-budget memory trained to retain frontier-specific predictive
    sufficient statistics;
    \item a stochastic revelation model used for information-, cost-, and
    risk-aware planning;
    \item a branch-complete HM3D protocol for evaluating every available
    frontier from an identical partial-observation state.
\end{enumerate}

\section{Workshop Scope}

The initial workshop study will focus on the benchmark, lineage-aware memory,
and three revelation outputs: occupancy, semantic features, and target
presence. The principal hypotheses are that (i) learned frontier memory
outperforms temporal and hand-designed caches under matched budgets,
(ii) lineage preservation matters most during split and merge events, and
(iii) lower counterfactual frontier regret translates into improved SPL.
Photorealistic future generation, large-scale cross-frontier factor graphs,
and real-robot online adaptation are reserved for the subsequent full study.

\begin{thebibliography}{10}
\footnotesize

\bibitem{yamauchi1997}
B.~Yamauchi.
\newblock A frontier-based approach for autonomous exploration.
\newblock In \textit{Proc.\ IEEE CIRA}, 1997.

\bibitem{frontiernet}
B.~Sun, H.~Chen, S.~Leutenegger, C.~Cadena, M.~Pollefeys, and H.~Blum.
\newblock FrontierNet: Learning visual cues to explore.
\newblock \textit{IEEE Robotics and Automation Letters}, 2025.

\bibitem{vlfm}
N.~Yokoyama, S.~Ha, D.~Batra, J.~Wang, and B.~Bucher.
\newblock VLFM: Vision-language frontier maps for zero-shot semantic
navigation.
\newblock In \textit{Proc.\ IEEE ICRA}, 2024.

\bibitem{ftm}
X.~Cui, Q.~Liu, Z.~Liu, and H.~Wang.
\newblock Frontier-enhanced topological memory with improved exploration
awareness for embodied visual navigation.
\newblock In \textit{Proc.\ ECCV}, 2024.

\bibitem{foresightnav}
H.~Shah, J.~Xing, N.~Messikommer, B.~Sun, M.~Pollefeys, and D.~Scaramuzza.
\newblock ForesightNav: Learning scene imagination for efficient exploration.
\newblock In \textit{CVPR Workshops}, 2025.

\bibitem{nwm}
A.~Bar, G.~Zhou, D.~Tran, T.~Darrell, and Y.~LeCun.
\newblock Navigation world models.
\newblock In \textit{Proc.\ CVPR}, 2025.

\bibitem{fomnav}
T.~Chabal, S.~Chen, J.~Ponce, and C.~Schmid.
\newblock FOM-Nav: Frontier-object maps for object goal navigation.
\newblock \textit{arXiv preprint arXiv:2512.01009}, 2025.

\bibitem{openfrontier}
E.~Padilla, B.~Sun, M.~Pollefeys, and H.~Blum.
\newblock OpenFrontier: General navigation with visual-language grounded
frontiers.
\newblock \textit{arXiv preprint arXiv:2603.05377}, 2026.

\end{thebibliography}

\clearpage
\appendix

\section{Implementation Progress Log}

Running record of implemented phases and measured results. Not part of the
submitted paper; kept here so every number in the main text can be traced to a
run that produced it.

\textbf{Phase 1 (complete).} Repository, conda environment, HM3D wiring, and a
one-command episode gate that saves RGB, depth, semantics, pose and an
occupancy map. Determinism verified: two runs at the same seed produce
identical action sequences, poses and map statistics. HM3D \emph{val} and
\emph{minival} scenes are available; \emph{train} scene meshes are not yet
downloaded, so Phase 5 generation is currently limited to val scenes.

\textbf{Phase 2 (complete).} Occupancy mapping from depth and ground-truth
pose; free--unknown boundary extraction, clustering, per-frontier centroid,
orientation and geometric information gain; navmesh planning; and the four
non-learning frontier policies. Table~\ref{tab:phase2} reports 50 episodes per
policy on HM3D-v2 val.

\begin{table}[h]
\centering
\footnotesize
\setlength{\tabcolsep}{3pt}
\caption{Phase 2 baseline frontier policies. 50 episodes per policy on an
identical episode list, 500-step budget. Intervals are 95\% confidence
half-widths. These are measured values, not placeholders; Tables~\ref{tab:ablations}
and~\ref{tab:lineage} remain to be filled.}
\label{tab:phase2}
\begin{tabular}{lcccc}
\toprule
Policy & SR \(\uparrow\) & SPL \(\uparrow\) & Det. & Area\\
       &                 &                  & rate & m\(^2\)\\
\midrule
Random               & 0.68\(\pm\)0.13 & 0.33\(\pm\)0.08 & 0.90 & 52.6\\
Nearest              & 0.72\(\pm\)0.13 & 0.40\(\pm\)0.08 & 0.96 & 47.7\\
Max information gain & 0.70\(\pm\)0.13 & 0.31\(\pm\)0.08 & 0.94 & 57.4\\
Gain \(-\) travel cost & 0.64\(\pm\)0.13 & 0.32\(\pm\)0.08 & 0.90 & 54.5\\
\bottomrule
\end{tabular}
\end{table}

\textbf{Conditions.} Three components use privileged simulator information and
must be labelled wherever these numbers appear. Frontier-crossing actions are
produced by habitat's navmesh planner, which is ground truth for the whole
scene including unobserved regions; stopping uses an oracle detector reading
ground-truth semantic instance ids; and success is scored at
\(1.0\,\mathrm{m}\) to the nearest goal viewpoint rather than habitat's
\(0.1\,\mathrm{m}\) default, so SR and SPL here are \emph{not} comparable to
published HM3D ObjectNav results.

\textbf{Reading of Phase 2.} The four policies are statistically
indistinguishable at \(n=50\): the confidence intervals overlap completely, so
the apparent ordering carries no signal. This is the correct baseline floor
for the paper's claim -- if prediction-based selection does not separate from
these, it has not earned its complexity. Two artefacts are worth carrying
forward. Nearest-frontier attains the best SPL while exploring the
\emph{least} area, because SPL rewards short paths and a loose success radius
lets most episodes resolve before exploration strategy matters; a tighter
radius should spread the policies apart. And gain-minus-cost is currently the
weakest policy, which indicates the travel-cost weight (\(1\,\mathrm{m^2}\)
per metre) is mistuned rather than that the rule is wrong -- it is the
baseline Eq.~\eqref{eq:planner} reduces to, so it deserves a sweep before it
is used as the reference.

\textbf{Known gaps.} Roughly \(15\%\) of episodes detect the target but never
call STOP: the approach plan fails and the agent resumes exploring. This is a
controller failure, not an exploration failure, and it caps SR for every
policy equally. Planner failures occur in 5--7 episodes per policy but are
always recovered by reselecting another frontier; none terminated an episode.

\textbf{Phase 3 (complete).} The canonical option
\(\Option_i=(\tau_i^{\mathrm{approach}},\tau_i^{\mathrm{cross}},H)\) is fixed
by frontier geometry and configuration alone, never by execution outcome, so
every candidate at a decision state is offered the same kind of chance. The
approach pose is the frontier's back-off point snapped to the navmesh and then
\emph{required to lie in free space the agent has observed}; the navmesh covers
the whole scene, so snapping alone can place an approach pose in territory
never seen, which would smuggle privileged geometry into the option definition
rather than merely into the controller. Probe distance is \(2.0\,\mathrm{m}\)
and \(H=12\) actions. Options that cannot be reached before crossing are
rejected with a recorded reason instead of silently thinning the branch set.

Branch rollouts drive the simulator agent directly rather than through the
episode stepper, so a counterfactual costs the real episode nothing and cannot
end it. A snapshot of the agent state and both maps is restored between
branches, and every branch is checked to begin from the recorded state.
\emph{Gate}: five decision states each executed three independent branches from
a bitwise-identical start, with state and maps verified restored and all
branches ending in distinct places.

\textbf{Phase 4 (complete).} Table~\ref{tab:phase4} reports the ground-truth
revelation over 27 branches from 7 decision states.
Figure~\ref{fig:branch} shows the verification view for one branch.

\begin{table}[h]
\centering
\footnotesize
\setlength{\tabcolsep}{3pt}
\caption{Ground-truth revelation, 27 branches from 7 decision states
(3.9 branches per state). Room category is derived, not ground truth.}
\label{tab:phase4}
\begin{tabular}{lccc}
\toprule
Revelation channel & mean & sd & max\\
\midrule
Newly revealed area \(\Delta M^{\mathrm{occ}}\) (m\(^2\)) & 8.63 & 5.21 & 22.78\\
\quad of which free space (m\(^2\)) & 6.33 & 4.29 & 19.19\\
Newly labelled semantics in \(\Delta M^{\mathrm{occ}}\) (m\(^2\)) & 7.44 & 5.11 & 21.45\\
New frontiers exposed & 2.70 & 2.16 & 8\\
Collisions during the option & 3.59 & 2.38 & 7\\
Collision rate (per action) & 0.11 & 0.09 & 0.35\\
Distance travelled (m) & 5.04 & 1.89 & 8.50\\
Actions executed & 39.0 & 10.6 & 64\\
\midrule
Crossings that reached unknown space & \multicolumn{3}{c}{24 / 27}\\
Branches where the target became visible & \multicolumn{3}{c}{8 / 27}\\
\bottomrule
\end{tabular}
\end{table}

\begin{figure*}[t]
\centering
\includegraphics[width=\textwidth]{figures/branch_example.png}
\caption{Verification view for one branch. Left to right, top row: the map
before crossing with the chosen frontier in red and the other candidates in
green; the map after, with \(\Delta M^{\mathrm{occ}}\) in blue; the executed
trajectory, solid for the approach and dashed for the crossing probe. Bottom
row: semantics newly labelled inside the revealed region; the egocentric view
at the decision state; and at the end of the crossing. Every quantity in the
revelation record is checkable against this figure, which is how the two
labelling defects below were found.}
\label{fig:branch}
\end{figure*}

\textbf{Counterfactual spread.} Branches from the same decision state differ
sharply in what they reveal -- within-state ratios between the best and worst
candidate range from \(1.6\times\) to \(115\times\)
(median \(3.1\times\)). This is the premise of the whole paper made
measurable: if every frontier revealed roughly the same information, there
would be nothing for a revelation model to predict and no regret to reduce.

\textbf{Two labelling defects found by manual verification.} Both would have
corrupted the prediction target silently.
\begin{itemize}
    \item \(\Delta M^{\mathrm{sem}}\) initially counted every cell that gained
    a semantic label, including cells already mapped geometrically but never
    labelled. That ran at \(3.26\times\) the correct figure in \(23/23\)
    records, and would have trained the model to predict re-labelling of space
    the agent had already seen. It is now restricted to labels gained inside
    \(\Delta M^{\mathrm{occ}}\).
    \item The room label was derived from every instance visible during a
    branch, including the room the agent started in. It is now derived only
    from instances the branch newly observed.
\end{itemize}

\textbf{Room category is the weakest channel.} HM3D-v0.2 ships no room-type
labels: every region returns a null category and a degenerate bounding box. The
label is therefore \emph{derived} from each region's object composition and is
recorded as derived in every stored example. Even after the fix above it skews
toward one class (kitchen in 13 of 27 branches), and spot checks disagree with
the egocentric view. Room prediction should either move to MP3D, which ships
genuine region categories, or be dropped from the reported channels; no number
in this paper should rest on it as it stands.

\end{document}